% CVPR 2026 Paper Template; see https://github.com/cvpr-org/author-kit

\documentclass[10pt,twocolumn,letterpaper]{article}

%%%%%%%%% PAPER TYPE  - PLEASE UPDATE FOR FINAL VERSION
% \usepackage{cvpr}              % To produce the CAMERA-READY version
\usepackage[review]{cvpr}      % To produce the REVIEW version
% \usepackage[pagenumbers]{cvpr} % To force page numbers, e.g. for an arXiv version

% Import additional packages in the preamble file, before hyperref
\input{preamble}

% It is strongly recommended to use hyperref, especially for the review version.
% hyperref with option pagebackref eases the reviewers' job.
% Please disable hyperref *only* if you encounter grave issues, 
% e.g. with the file validation for the camera-ready version.
%
% If you comment hyperref and then uncomment it, you should delete *.aux before re-running LaTeX.
% (Or just hit 'q' on the first LaTeX run, let it finish, and you should be clear).
\definecolor{cvprblue}{rgb}{0.21,0.49,0.74}
\usepackage[pagebackref,breaklinks,colorlinks,allcolors=cvprblue]{hyperref}

%%%%%%%%% PAPER ID  - PLEASE UPDATE
\def\paperID{*****} % *** Enter the Paper ID here
\def\confName{CVPR}
\def\confYear{2026}

%%%%%%%%% TITLE - PLEASE UPDATE
\title{Weight-Space Prediction for Federated Continual Learning: \\
A Systematic Study of 39 Loss Functions via Topology and Optimal Transport}

%%%%%%%%% AUTHORS - PLEASE UPDATE
\author{Anonymous Author(s)\\
Anonymous Institution\\
{\tt\small anonymous@institution.edu}
}

\begin{document}
\maketitle
\begin{abstract}
Federated Continual Learning (FCL) requires composing knowledge from separately trained models without access to raw data---a challenge fundamentally tied to the geometry of neural network weight spaces. We introduce a \emph{hyper-representation} framework that learns to predict merged model weights directly from paired individual-task weight vectors, enabling zero-shot federated composition. A 4.02M-parameter dual-encoder TransformerAE is trained to map $(w_{\tau_1}, w_{\tau_2}) \mapsto w_{\tau_1 \cup \tau_2}$ across 39 loss functions spanning reconstruction, spectral, optimal transport, information-theoretic, and topological families. We evaluate 56 completed tiny-scale experiments across three data-overlap scenarios, achieving mean post-finetune in-distribution (ID) CNN accuracy of 80.6\%, 76.6\%, and 84.5\% for overlaps 0, 1, and 2 respectively, starting from near-random $\sim$14\% initialization---representing a $\mathbf{+66}$ percentage-point improvement after only 5 fine-tuning epochs on 100 samples. Loss function choice produces systematic effects: log-norm regularization dominates under disjoint data (overlap 0), spectral losses (FFT+MelSpec) peak at partial overlap, and optimal transport (Sinkhorn) excels at high overlap. Multiparameter persistent homology and Gromov--Wasserstein analysis reveal that OT-based losses induce topologically simpler weight spaces (GW $\approx$ 7--20) compared to reconstruction losses (GW $\approx$ 32), correlating with transfer efficiency. \textcolor{gray}{[Topology, HMMR segmentation, and full tournament results (small/medium/large): TBD.]}
\end{abstract}    
\section{Introduction}
\label{sec:intro}

\textbf{Federated Continual Learning (FCL)} poses a fundamental challenge: given models independently trained on disjoint task subsets $\tau_1, \tau_2 \subset \mathcal{T}$, can we produce a composed model for $\tau_1 \cup \tau_2$ \emph{without access to raw training data}? Classical federated averaging~\cite{mcmahanCommunicationEfficientLearningDeep2017} requires simultaneous client participation and fails under continual shift; gradient-based merging~\cite{kirkpatrickOvercomingCatastrophicForgetting2017} requires data access. We address this via \emph{weight-space prediction}: learn a mapping $f: (w_{\tau_1}, w_{\tau_2}) \mapsto \hat{w}_{\tau_1 \cup \tau_2}$ from paired weight vectors, treating model merging as a \emph{regression problem in weight space}.

This framing has a critical advantage: at inference time, no data, gradients, or communication overhead is needed---the merged model is obtained in a single forward pass through $f$, followed by lightweight fine-tuning ($\leq$5 epochs, 100 samples) to restore task-specific calibration.

\textbf{The challenge of loss function design.} The weight prediction problem spans a high-dimensional, geometrically complex manifold ($d = 2{,}464$ parameters for our CNN). The choice of training objective for $f$ profoundly shapes the learned weight space geometry: does the predicted $\hat{w}$ preserve distributional properties (motivating optimal transport~\cite{villaniOptimalTransportOld2009}), spectral structure (motivating FFT-based losses), or topological invariants (motivating persistence-based losses~\cite{hoferDeepLearningTopological2019})? No prior work has systematically characterized this trade-off.

\textbf{Our approach.} We train a 4.02M-parameter dual-encoder \textbf{TransformerAE} on a \emph{weight zoo}---a database of 4,052 CNN weight vectors organized by task label, epoch, and activation---and evaluate 39 distinct loss functions spanning 7 families across three data-overlap regimes (disjoint, partial, high). Each trained model is evaluated by initializing a CNN with predicted weights and measuring accuracy after 5-epoch fine-tuning on 100 held-out samples, providing a directly interpretable downstream metric.

\subsection{Contributions}

\begin{enumerate}
    \item \textbf{FCL weight prediction framework}: A dual-encoder TransformerAE that maps paired task weight vectors to merged-model weights, achieving $>$79\% CNN accuracy after 5-epoch fine-tuning from a 14\% random baseline---across all overlap conditions and loss families (\S\ref{sec:results}).

    \item \textbf{Systematic 39-loss tournament}: The first large-scale comparison of reconstruction, spectral, optimal-transport, information-theoretic, geometric, autoregressive, and topological losses for weight-space prediction, revealing overlap-dependent winner profiles (\S\ref{sec:results}).

    \item \textbf{Topological analysis of loss-induced weight geometry}: Using Gromov--Wasserstein distance, multiparameter persistent homology, and HMMR time-series segmentation, we characterize how loss choice shapes topological structure. OT losses yield lower GW distances ($7$--$20$) vs.\ reconstruction losses ($\sim$32), correlating with transfer efficiency (\S\ref{sec:results}). \textcolor{gray}{[Full topology results: TBD.]}

    \item \textbf{Multi-stage tournament protocol}: A progressive elimination design (tiny $\to$ small $\to$ medium $\to$ large model) that identifies robust loss-overlap pairings under increasing representational capacity. \textcolor{gray}{[Small/medium/large stages: TBD.]}
\end{enumerate}

\subsection{Why Loss Function Matters for Weight-Space Prediction}

The loss $\mathcal{L}(f(w_{\tau_1}, w_{\tau_2}),\, w_{\tau_1\cup\tau_2})$ determines the geometry of the learned prediction manifold. We identify seven families with distinct inductive biases:

\begin{itemize}
    \item \textbf{Reconstruction} (MSE, MAE, MAPE, Quantile): Direct $\ell_p$ distances; establish geometric baselines.
    \item \textbf{Spectral} (FFT, MelSpec, LW\_FFT): Frequency-domain alignment; sensitive to oscillatory weight patterns.
    \item \textbf{Optimal transport} (Sinkhorn, LW\_Sinkhorn, Sinkhorn+KL): Wasserstein-type matching; preserves distributional shape.
    \item \textbf{Information-theoretic} (KL, JS): Probabilistic divergences; promote distributional coverage.
    \item \textbf{Geometric/norm} (Frobenius, LogNorm, FIM): Matrix conditioning and scale control.
    \item \textbf{Autoregressive} (AUTO): Temporal consistency of weight sequences.
    \item \textbf{Topological} (DiffPers, RTD, PersLandscape, and layerwise variants): Directly optimize persistent homology features of predicted weight manifolds. \textcolor{gray}{[Results: TBD.]}
\end{itemize}

The remainder of this paper is organized as follows: \S\ref{sec:related} reviews related work. \S\ref{sec:method} details the weight-zoo construction, TransformerAE architecture, and loss taxonomy. \S\ref{sec:experiments} describes the experimental protocol. \S\ref{sec:results} presents current and forthcoming results. \S\ref{sec:conclusion} concludes.

\section{Related Work}
\label{sec:related}

\subsection{Topological Data Analysis for Neural Networks}

Topological data analysis (TDA) has emerged as a powerful framework for understanding neural network behavior. Carlsson and Gabrielsson~\cite{carlssonTopologicalApproachesDeep2018} pioneered the application of persistent homology to deep learning, demonstrating that topological features of activation spaces correlate with network performance. Ballester \etal~\cite{ballesterTopologicalDataAnalysis2024} provide a comprehensive survey of TDA methods for neural network analysis, covering applications from architecture design to interpretability.

Recent work has focused on using persistence diagrams to analyze network stability~\cite{akaiExperimentalStabilityAnalysis2021} and generalization~\cite{birdalIntrinsicDimensionPersistent2021}. Birdal \etal show that intrinsic dimension and persistent homology jointly predict generalization, establishing a theoretical connection between topology and learning dynamics. Our work extends these insights to weight space, where topological structure directly reflects optimization trajectories.

\subsection{Multiparameter Persistent Homology}

Classical persistent homology tracks topological features along a single filtration parameter. Carlsson and Zomorodian~\cite{carlssonTheoryMultidimensionalPersistence2009} introduced the theory of multidimensional persistence, showing that no complete discrete invariant exists for multifiltrations. Despite this theoretical limitation, recent advances have made multiparameter persistence practical. Botnan \etal~\cite{botnanBottleneckStabilityRank2024} establish bottleneck stability for rank decompositions, providing theoretical foundations for robust multiparameter analysis.

Carrière and Blumberg~\cite{carriereMultiparameterPersistenceImages} introduce multiparameter persistence images for machine learning, demonstrating that vectorized representations enable statistical analysis while preserving topological information. We build on this work, applying multiparameter persistence to neural network weight spaces where multiple training parameters naturally define a multifiltration.

\subsection{Neural Network Weight Space Analysis}

Understanding weight space geometry is crucial for optimization and generalization. The Neural Tangent Kernel (NTK) framework~\cite{adlamNeuralTangentKernel2020} provides theoretical insights into training dynamics in the infinite-width limit. However, practical networks exhibit rich finite-width phenomena that NTK theory cannot capture.

Recent work on representation topology~\cite{barannikovRepresentationTopologyDivergence2022} introduces divergence measures based on topological features, enabling comparison of learned representations. Our hyper-representation framework extends this idea to weight spaces, treating entire networks as points in a learned embedding space.

\subsection{Loss Function Design}

Loss function choice profoundly affects optimization landscapes and learned representations. Beyond standard reconstruction losses, recent work explores spectral losses for frequency-aware training, optimal transport losses for distributional matching~\cite{230909442UseKantorovichRubinstein2025}, and information-theoretic objectives for representation learning. Our systematic study of 18 loss functions reveals how different objectives shape weight space topology, with regularized combinations exhibiting superior stability.

\subsection{Distributed Computation of Persistence}

Computing persistent homology for large-scale neural networks requires efficient algorithms. Bauer \etal~\cite{bauerDistributedComputationPersistent2013} introduce distributed methods for persistence computation, enabling analysis of high-dimensional datasets. We leverage these techniques to analyze weight spaces with thousands of dimensions, computing multiparameter persistence diagrams for entire training trajectories.

\section{Methodology}
\label{sec:method}

\subsection{Problem Formulation: Weight-Space Prediction for FCL}

Let $\mathcal{W} = \mathbb{R}^d$ be the weight space of a CNN architecture ($d = 2{,}464$). For task subsets $\tau_1, \tau_2 \subset \mathcal{T}$ (MNIST classes), let $w_\tau \in \mathcal{W}$ denote the weights of a CNN trained to convergence on $\tau$. Our goal is to learn a \emph{merging predictor}:
\begin{equation}
f_\theta: \mathcal{W} \times \mathcal{W} \to \mathcal{W}, \quad
f_\theta(w_{\tau_1}, w_{\tau_2}) \approx w_{\tau_1 \cup \tau_2}
\label{eq:merge_objective}
\end{equation}
trained on a database of $(w_{\tau_1}, w_{\tau_2}, w_{\tau_1\cup\tau_2})$ triplets. The predicted $\hat{w} = f_\theta(w_{\tau_1}, w_{\tau_2})$ is evaluated by fine-tuning a CNN initialized at $\hat{w}$ on a fixed 100-sample held-out set for 5 epochs, then measuring in-distribution (ID) classification accuracy on $\tau_1 \cup \tau_2$.

\subsection{Weight Zoo Construction}

We build a \emph{weight zoo}: 4,052 CNN weight vectors trained on all valid task subsets of MNIST classes $\{0,\ldots,9\}$, indexed by $(label\_set, activation, epoch)$. The lookup index enables $O(1)$ retrieval per pair. Training data statistics depend on the overlap regime:

\begin{table}[h]
\centering\small
\begin{tabular}{lccc}
\hline
\textbf{Overlap} & \textbf{Train pairs} & \textbf{Val pairs} & \textbf{Test pairs} \\
\hline
0 (disjoint)  & 32,812  & 7,030  & 7,030 \\
1 (partial)   & 130,620 & 27,990 & 27,990 \\
2 (high)      & 130,620 & 27,990 & 27,990 \\
\hline
\end{tabular}
\caption{Dataset sizes per overlap regime. Overlap controls how many valid $(w_{\tau_1}, w_{\tau_2}, w_{\tau_1\cup\tau_2})$ triplets exist.}
\label{tab:dataset}
\end{table}

\textbf{Layer-wise normalization.} Weights are standardized per layer segment using boundaries $[0, 208, 1{,}414, 1{,}514, 2{,}254, 2{,}464]$ corresponding to (conv1\_w, conv1\_b, conv2\_w, conv2\_b, fc\_w, fc\_b). This prevents dominant layers from dominating the training signal.

\subsection{TransformerAE Architecture}

The dual-encoder TransformerAE processes both input weight vectors in parallel:
\begin{equation}
\hat{w} = \text{Dec}\!\left(\text{Enc}(w_{\tau_1}) \oplus \text{Enc}(w_{\tau_2})\right)
\label{eq:dual_enc}
\end{equation}
where $\oplus$ denotes concatenation-then-projection. The encoder treats weights as sequences of 16-parameter patches (154 patches for $d=2{,}464$):

\begin{itemize}
    \item \textbf{Encoder}: 3 transformer layers, 2 attention heads, dropout $p=0.1$; processes each input independently
    \item \textbf{Bottleneck}: latent merge of dual encodings via linear projection
    \item \textbf{Decoder}: 2 transformer layers for weight reconstruction
    \item \textbf{Total parameters}: 4.02M
\end{itemize}

Attention is computed as:
\begin{equation}
\text{Attn}(Q,K,V) = \text{softmax}\!\left(\tfrac{QK^\top}{\sqrt{d_k}}\right)V
\label{eq:attention}
\end{equation}

\subsection{Loss Function Taxonomy (39 functions)}

All loss functions compare the predicted weights $\hat{w}$ to ground-truth merged weights $w^*$. We organize them into three hierarchical levels:

\textbf{Level 1---Individual losses (13 base):}
\begin{itemize}
    \item \textit{Reconstruction}: MSE, MAE, MAPE, Quantile
    \item \textit{Spectral}: FFT, MelSpec
    \item \textit{Optimal transport}: Sinkhorn (Wasserstein-2 via \texttt{geomloss})
    \item \textit{Information-theoretic}: KL, JS
    \item \textit{Geometric/norm}: Frobenius, LogNorm, FIM
    \item \textit{Autoregressive}: AUTO
\end{itemize}

\textbf{Level 2---Layer-wise variants (LW\_*):}
Each base loss is applied independently to each of the 5 CNN layers, then summed:
\begin{equation}
\mathcal{L}_{\text{LW}} = \sum_{l=1}^{5} \mathcal{L}(w^{(l)}_{\text{pred}},\, w^{(l)}_*)
\label{eq:layerwise}
\end{equation}
Boundaries: $[0, 208, 1{,}414, 1{,}514, 2{,}254, 2{,}464]$.

\textbf{Level 3---Regularized combinations (13 combos):}
\begin{equation}
\mathcal{L}_{\text{reg}} = \mathcal{L}_{\text{main}}(\hat{w}, w^*) + \lambda\,\mathcal{L}_{\text{aux}}(\hat{w}, w^*)
\label{eq:regularized}
\end{equation}
Selected pairs and weights: MSE+$0.05\times$Frobenius, MSE+$0.1\times$LogNorm, FFT+$0.1\times$MelSpec, Sinkhorn+$0.15\times$KL, and layerwise equivalents. The regularizer $\lambda$ is chosen to keep the auxiliary term sub-dominant ($\lambda \leq 0.15$).

\textbf{Topological losses (6, Level 3):} DiffPers, RTD, LW\_DiffPers, LW\_RTD, MSE+$0.1\times$DiffPers, LW\_MSE+$0.01\times$LW\_PersLandscape. These directly optimize features of the Vietoris-Rips persistence diagram of the predicted weight manifold. \textcolor{gray}{[Results pending.]}

\subsection{Topological Analysis}

\textbf{Gromov--Wasserstein distance.} For each trained model we compute $d_{\text{GW}}$ between the weight distribution of predicted outputs and ground-truth weights, measuring geometric dissimilarity beyond pointwise MSE.

\textbf{Multiparameter persistent homology.} We construct a bifiltration over (weight magnitude, layer position) and compute rank invariants on a $10\times10$ grid using \texttt{multipers}~\cite{lesnickInteractiveVisualization2015}. Features extracted: $\beta_0, \beta_1$, persistence entropy $E = -\sum_i p_i\log p_i$, effective rank $r_{\text{eff}} = e^E$, total persistence.

\textbf{HMMR time-series segmentation.} Training loss curves are segmented using changepoint detection to identify distinct training phases (fast descent, plateau, oscillation). \textcolor{gray}{[Results pending.]}

\subsection{Evaluation Protocol}

For each trained $f_\theta$:
\begin{enumerate}
    \item Run $f_\theta$ on all test-set pairs $(w_{\tau_1}, w_{\tau_2})$ to obtain $\hat{w}$.
    \item Initialize CNN at $\hat{w}$; record initial accuracy (pre-finetune).
    \item Fine-tune for 5 epochs on 100 fixed samples drawn with seed 42 (identical across all experiments).
    \item Record final in-distribution (ID) and out-of-distribution (OOD) accuracy.
    \item Compute reconstruction metrics: cosine similarity, Wasserstein-1, GW distance.
\end{enumerate}

All 56 completed experiments share \emph{identical} evaluation samples (seed 42, $n=100$), ensuring strict comparability.

\section{Experimental Setup}
\label{sec:experiments}

\subsection{Base CNN Architecture}

The target CNN (whose weights we predict) has the following structure:
\begin{itemize}
    \item Conv1: $1\to32$ filters, $3\times3$, ReLU
    \item Conv2: $32\to64$ filters, $3\times3$, ReLU
    \item FC1: $64\times9\times9 \to 128$, LeakyReLU
    \item FC2: $128 \to C$ (C = number of task classes, variable)
    \item \textbf{Total fixed weights}: $d = 2{,}464$ parameters
    \item \textbf{Layer boundaries}: $[208, 1{,}414, 1{,}514, 2{,}254, 2{,}464]$
\end{itemize}

\subsection{TransformerAE Training Protocol}

\begin{itemize}
    \item \textbf{Epochs}: 200 (all experiments converge; checkpointed at best val loss)
    \item \textbf{Batch size}: 24
    \item \textbf{Optimizer}: AdamW, lr $= 10^{-4}$, weight decay $10^{-5}$
    \item \textbf{Scheduler}: CosineAnnealingLR, $T_{\max}=200$
    \item \textbf{Dropout}: 0.1 in encoder and decoder
    \item \textbf{Checkpoint frequency}: Best val-loss + final epoch
\end{itemize}

\subsection{Tournament Design}

We run a progressive multi-stage tournament. \textbf{Stage 1 (tiny)} uses 39 curated loss functions across overlaps $\{0,1,2\}$ (117 total experiments). Subsequent stages use model-size progression with loss selection based on prior-stage rankings:

\begin{table}[h]
\centering\small
\begin{tabular}{lcccc}
\hline
\textbf{Stage} & \textbf{Model} & \textbf{Selection} & \textbf{Exps.} & \textbf{Status} \\
\hline
1 & tiny   & All 39 losses     & 117 & 56 done, 63 running \\
2 & small  & top 30\% + bot 10\% & $\sim$48 & \textcolor{gray}{TBD} \\
3 & medium & top 20\% + bot 5\%  & $\sim$27 & \textcolor{gray}{TBD} \\
4 & large  & top 10\% + bot 2.5\% & $\sim$15 & \textcolor{gray}{TBD} \\
\hline
\end{tabular}
\caption{Tournament schedule. ``bot'' rows test whether worst performers are loss-specific failures or architecture-limited.}
\label{tab:tournament}
\end{table}

The ``bottom carrier'' rule---advancing worst performers alongside top---tests whether low-overlap failures recover under larger model capacity, separating architecture-limited from loss-limited behavior.

\subsection{Completed Experiments (Stage 1, Partial)}

Of the 117 planned Stage-1 experiments, \textbf{56 have completed}: 20 for overlap 0, 18 each for overlaps 1 and 2. These cover the following loss families:
\begin{itemize}
    \item Reconstruction: MSE, MAE, MAPE (+ layerwise)
    \item Spectral: FFT, FFT+$0.1\times$MelSpec (+ LW variants)
    \item Optimal transport: Sinkhorn, LW\_Sinkhorn
    \item Norm-regularized: MSE+$0.05\times$Frobenius, MSE+$0.1\times$LogNorm (+ LW)
    \item MAPE regularized: MAPE+$0.1\times$JS, LW\_MAPE+$0.1\times$LW\_JS
    \item Autoregressive: AUTO
    \item LWLN (layerwise log-norm normalization)
\end{itemize}

Running ($n=63$): DiffPers, RTD, LW\_DiffPers, LW\_RTD, persistence combos, MelSpec, LW\_FFT, Sinkhorn+KL, FIM, Quantile, and variants.

\subsection{CNN Fine-Tuning Evaluation}

For each checkpoint (epochs 25, 50, 75, 100, 125, 150, 175, 200), we:
\begin{enumerate}
    \item Sample 10 test pairs from the held-out set
    \item Run inference to obtain 10 predicted weight vectors
    \item Initialize 10 CNNs and fine-tune each for 5 epochs (batch size 24) on 100 samples
    \item Record ID accuracy (on $\tau_1\cup\tau_2$ classes) and OOD accuracy (on held-out classes)
\end{enumerate}

\textbf{Fair evaluation:} The final comparative ranking uses a \emph{fixed} 100-sample set (seed 42, identical across all experiments) and the best-epoch checkpoint. This ensures all 39 losses are compared under identical conditions.

\subsection{Computational Resources}

All experiments run on a single \textbf{NVIDIA RTX 5060 Ti (16.61 GB)} GPU. Wall-clock times per experiment:
\begin{itemize}
    \item Overlap 0: $\approx$45--90 min/experiment (32K pairs)
    \item Overlap 1/2: $\approx$2--4 h/experiment (130K pairs)
    \item Topological losses (DiffPers, RTD): $\approx$2$\times$ slower due to persistent homology computation
\end{itemize}

Total Stage-1 budget: $\approx 117 \times 2.5\text{h avg} \approx 293$ GPU-hours.

\section{Results and Analysis}
\label{sec:results}

\noindent\textbf{Status note.} This section reports results from \textbf{56 completed} Stage-1 tiny experiments (baselines) with 63 topology/spectral experiments still running. All accuracy figures use CNN fine-tuning on the fixed 100-sample seed-42 test set. Rows marked \textcolor{gray}{[TBD]} will be filled upon completion.

\subsection{CNN Fine-Tuning Performance (56 Completed Experiments)}
\label{sec:res_cnn}

Table~\ref{tab:cnn_main} reports mean ID accuracy after 5-epoch fine-tuning from predicted weights. All 56 experiments use identical evaluation conditions (seed 42, $n{=}100$ samples).

\begin{table}[t]
\centering\small
\setlength{\tabcolsep}{4pt}
\begin{tabular}{lcccc}
\hline
\textbf{Loss} & \textbf{OV0} & \textbf{OV1} & \textbf{OV2} & \textbf{Avg} \\
\hline
\multicolumn{5}{l}{\textit{--- Reconstruction ---}} \\
MSE           & 79.3 & 76.8 & 84.4 & 80.2 \\
MAE           & 81.2 & 76.2 & 84.8 & 80.7 \\
MAPE          & 80.2 & 75.8 & 84.1 & 80.0 \\
LW\_MSE       & 80.0 & 75.9 & 84.8 & 80.2 \\
LW\_MAE       & 81.3 & 77.1 & 84.8 & 81.1 \\
LW\_MAPE      & 80.0 & 75.7 & 84.1 & 79.9 \\
\hline
\multicolumn{5}{l}{\textit{--- Norm-regularized ---}} \\
MSE+$0.05\times$Frob  & 79.6 & 77.2 & 84.4 & 80.4 \\
MSE+$0.1\times$LogN   & \textbf{82.4} & 76.9 & \textbf{84.8} & \textbf{81.4} \\
LW\_MSE+$0.05\times$LW\_Frob & 79.7 & \textbf{77.6} & 84.6 & 80.7 \\
LW\_MSE+$0.1\times$LW\_LogN  & 82.1 & 76.9 & 84.8 & 81.3 \\
MAPE+$0.1\times$JS    & 79.8 & 75.5 & 83.9 & 79.7 \\
LW\_MAPE+$0.1\times$LW\_JS & 79.8 & 75.2 & 83.8 & 79.6 \\
\hline
\multicolumn{5}{l}{\textit{--- Spectral ---}} \\
FFT           & 80.2 & 77.2 & 84.5 & 80.6 \\
FFT+$0.1\times$MelSpec & 79.4 & \textbf{77.7} & 84.4 & 80.5 \\
\hline
\multicolumn{5}{l}{\textit{--- Optimal transport ---}} \\
Sinkhorn      & 81.4 & 76.7 & 84.8 & 81.0 \\
LW\_Sinkhorn  & 81.6 & 76.9 & \textbf{85.1} & \textbf{81.2} \\
\hline
\multicolumn{5}{l}{\textit{--- Other ---}} \\
AUTO          & 81.5 & 76.6 & 84.7 & 80.9 \\
LWLN          & 81.3 & 76.4 & 84.6 & 80.8 \\
\hline
\multicolumn{5}{l}{\textit{--- Topology / Spectral (new) ---}} \\
DiffPers      & \textcolor{gray}{TBD} & \textcolor{gray}{TBD} & \textcolor{gray}{TBD} & \textcolor{gray}{TBD} \\
RTD           & \textcolor{gray}{TBD} & \textcolor{gray}{TBD} & \textcolor{gray}{TBD} & \textcolor{gray}{TBD} \\
MSE+DiffPers  & \textcolor{gray}{TBD} & \textcolor{gray}{TBD} & \textcolor{gray}{TBD} & \textcolor{gray}{TBD} \\
LW\_MSE+LW\_PersL. & \textcolor{gray}{TBD} & \textcolor{gray}{TBD} & \textcolor{gray}{TBD} & \textcolor{gray}{TBD} \\
MelSpec       & \textcolor{gray}{TBD} & \textcolor{gray}{TBD} & \textcolor{gray}{TBD} & \textcolor{gray}{TBD} \\
LW\_FFT       & \textcolor{gray}{TBD} & \textcolor{gray}{TBD} & \textcolor{gray}{TBD} & \textcolor{gray}{TBD} \\
Sinkhorn+KL   & \textcolor{gray}{TBD} & \textcolor{gray}{TBD} & \textcolor{gray}{TBD} & \textcolor{gray}{TBD} \\
Quantile      & \textcolor{gray}{TBD} & \textcolor{gray}{TBD} & \textcolor{gray}{TBD} & \textcolor{gray}{TBD} \\
FIM           & \textcolor{gray}{TBD} & \textcolor{gray}{TBD} & \textcolor{gray}{TBD} & \textcolor{gray}{TBD} \\
\hline
\textbf{Mean (done)} & \textbf{80.6} & \textbf{76.6} & \textbf{84.5} & \textbf{80.6} \\
\textbf{Std (done)}  & 1.0  & 0.7  & 0.4  & -- \\
\textbf{Init (pre-FT)} & 13.3 & 14.5 & 13.8 & 13.9 \\
\hline
\end{tabular}
\caption{\textbf{Stage-1 (tiny model) CNN fine-tuning accuracy (\%).} OV = overlap regime. \textbf{Bold}: best per column among completed runs. Initial accuracy ($\sim$14\%) reflects near-random CNN initialization from predicted weights. All 56 models trained 200 epochs; evaluated at best-val checkpoint.}
\label{tab:cnn_main}
\end{table}

\textbf{Key observations:}
\begin{itemize}
    \item \textbf{Near-uniform floor}: All 56 completed losses achieve $\geq$75\% ID accuracy, confirming weight-space prediction as a viable FCL strategy regardless of loss family.
    \item \textbf{Overlap-dependent winner}: MSE+$0.1\times$LogNorm leads OV0 (82.4\%); FFT+$0.1\times$MelSpec leads OV1 (77.7\%); LW\_Sinkhorn leads OV2 (85.1\%). No single loss dominates all three regimes.
    \item \textbf{OT robustness}: Sinkhorn and LW\_Sinkhorn rank top-3 in two of three overlaps, exhibiting the lowest cross-overlap variance ($\sigma = 0.4\%$ for LW\_Sinkhorn vs $1.8\%$ for spectral losses).
    \item \textbf{Layerwise versions help}: For norm-regularized and MAPE losses, LW variants match or exceed global counterparts (+0.0 to +0.3\%).
    \item \textbf{+66 pp gain}: Mean improvement over initial accuracy across all 56 experiments: $+66.7$ percentage points.
\end{itemize}

\subsection{Gromov--Wasserstein Geometric Analysis}

The GW distance $d_{\text{GW}}$ measures structural dissimilarity between the distribution of predicted weights and ground-truth weights, independent of pointwise alignment.

\begin{table}[h]
\centering\small
\begin{tabular}{lcccc}
\hline
\textbf{Loss family} & \textbf{OV0} & \textbf{OV1} & \textbf{OV2} & \textbf{Mean} \\
\hline
MAE / LW\_MAE       & 32.0 & 22.7 & 21.3 & 25.3 \\
MSE / LW\_MSE       & 31.9 & 22.7 & 22.1 & 25.6 \\
LogNorm combos       & 32.2 & 31.5 & 30.1 & 31.3 \\
FFT / FFT+MelSpec    & 32.0 & 17.8 & 17.9 & 22.6 \\
AUTO                 & 32.1 & 23.8 & 13.6 & 23.2 \\
\textbf{Sinkhorn}   & 23.1 & \textbf{6.6}  & 16.7 & 15.5 \\
\textbf{LW\_Sinkhorn} & 19.7 & 10.2 & \textbf{10.4} & \textbf{13.4} \\
LWLN                 & 32.2 & 21.3 & 15.8 & 23.1 \\
\hline
\end{tabular}
\caption{\textbf{Gromov--Wasserstein distance by loss family (lower = better geometric match).} OT-based losses (Sinkhorn) achieve dramatically lower GW distances, indicating they produce weight distributions geometrically closer to ground truth---even when CNN accuracy differences are modest.}
\label{tab:gw}
\end{table}

OT losses (Sinkhorn, LW\_Sinkhorn) achieve GW $\approx 6$--$20$ vs.\ reconstruction losses ($\approx 22$--$32$). This suggests OT-trained models produce weight manifolds that are geometrically closer to ground truth at a distributional level, not just pointwise---potentially explaining their superior robustness at high overlap.

\subsection{Training Dynamics}

\begin{table}[h]
\centering\small
\begin{tabular}{lcc}
\hline
\textbf{Loss family} & \textbf{Best val loss (OV0)} & \textbf{Convergence pattern} \\
\hline
MAE / LW\_MAE  & 0.807 / 0.818 & Smooth, monotone \\
MSE / LW\_MSE  & 1.015 / 1.014 & Smooth, monotone \\
LogNorm combos & 1.256 / 1.237 & Slower initial, better generalization \\
AUTO           & 14.10          & Different scale; fast plateau \\
MAPE variants  & $\approx$100   & \% scale; fast convergence \\
Sinkhorn       & 1265.9         & Wasserstein scale; high variance \\
LW\_Sinkhorn   & 253.2          & Wasserstein scale; stable \\
FFT variants   & $\approx$2500  & Spectral power scale; monotone \\
\hline
\end{tabular}
\caption{\textbf{Training dynamics summary (Overlap 0).} Loss values are not comparable across families (different units). All experiments reach convergence within 200 epochs.}
\label{tab:dynamics}
\end{table}

\textit{Note:} Loss values in Table~\ref{tab:dynamics} are in family-specific units and not directly comparable. All experiments run the full 200 epochs with cosine decay; no early-stopping failures observed.

\subsection{Topological Analysis \textcolor{gray}{[Partial]}}

\textbf{GW topology} (reported above, Table~\ref{tab:gw}) is the only topology metric currently computed for all 56 experiments. Full GUDHI persistence, multiparameter analysis (HMMR segmentation), and DiffPers/RTD loss-effect analysis are pending completion of the 63 running experiments.

\begin{table}[h]
\centering\small
\begin{tabular}{lcccc}
\hline
\textbf{Metric} & \textbf{OV0} & \textbf{OV1} & \textbf{OV2} & \textbf{Source} \\
\hline
GW distance (mean)    & 29.4 & 19.8 & 18.1 & \checkmark{} Done \\
Betti numbers $\beta_0,\beta_1$ & \textcolor{gray}{TBD} & \textcolor{gray}{TBD} & \textcolor{gray}{TBD} & Running \\
Persistence entropy   & \textcolor{gray}{TBD} & \textcolor{gray}{TBD} & \textcolor{gray}{TBD} & Running \\
Effective rank $r_{\text{eff}}$ & \textcolor{gray}{TBD} & \textcolor{gray}{TBD} & \textcolor{gray}{TBD} & Running \\
HMMR segments/exp.    & \textcolor{gray}{TBD} & \textcolor{gray}{TBD} & \textcolor{gray}{TBD} & Running \\
DiffPers vs.\ MSE $\Delta$acc & \textcolor{gray}{TBD} & \textcolor{gray}{TBD} & \textcolor{gray}{TBD} & Running \\
LW\_PersLandscape $\Delta$acc & \textcolor{gray}{TBD} & \textcolor{gray}{TBD} & \textcolor{gray}{TBD} & Running \\
\hline
\end{tabular}
\caption{\textbf{Topological analysis status.} GW distances are confirmed from 56 completed experiments. All other topology metrics await completion of the 63 running experiments.}
\label{tab:topo_status}
\end{table}

\subsection{Tournament Progression \textcolor{gray}{[TBD]}}

\begin{table}[h]
\centering\small
\begin{tabular}{lccccc}
\hline
\textbf{Stage} & \textbf{Model} & \textbf{$n$ losses} & \textbf{Best acc} & \textbf{Worst acc} & \textbf{Status} \\
\hline
1 & tiny   & 39 & \textcolor{gray}{TBD} & \textcolor{gray}{TBD} & 56 done \\
2 & small  & $\sim$16 & \textcolor{gray}{NaN} & \textcolor{gray}{NaN} & \textcolor{gray}{TBD} \\
3 & medium & $\sim$9  & \textcolor{gray}{NaN} & \textcolor{gray}{NaN} & \textcolor{gray}{TBD} \\
4 & large  & $\sim$5  & \textcolor{gray}{NaN} & \textcolor{gray}{NaN} & \textcolor{gray}{TBD} \\
\hline
\end{tabular}
\caption{\textbf{Multi-stage tournament results.} Stage 1 in progress; stages 2--4 pending. Best/worst accuracy values will be filled after each stage's fair evaluation run.}
\label{tab:tournament_results}
\end{table}

\subsection{Key Findings (Current)}

\begin{enumerate}
    \item \textbf{Universal viability}: All 18 baseline loss families tested achieve $\geq$75\% CNN fine-tuning accuracy from predicted weights, confirming weight-space FCL prediction is robust to loss choice.

    \item \textbf{Overlap-dependent winner profile}: LogNorm regularization wins disjoint data (OV0: 82.4\%); spectral FFT+MelSpec wins partial overlap (OV1: 77.7\%); OT (LW\_Sinkhorn) wins high overlap (OV2: 85.1\%). This strongly motivates overlap-aware loss selection.

    \item \textbf{OT produces geometrically closer weights}: Sinkhorn-family losses reduce GW distance by $47$--$79\%$ vs.\ MSE/MAE baselines (Table~\ref{tab:gw}), suggesting distributional shape preservation is a key inductive bias for weight merging.

    \item \textbf{Layerwise normalization is reliably beneficial}: LW variants match or exceed global counterparts in all 6 family comparisons; LWLN (layerwise log-norm) achieves 81.3\% avg vs.\ 80.6\% for standard MSE.

    \item \textbf{Topology/persistence results pending}: Whether DiffPers, RTD, and PersLandscape regularization improve upon these baselines is the core open question; results expected within the tournament window.
\end{enumerate}

\section{Conclusion}
\label{sec:conclusion}

We presented a \textbf{weight-space prediction framework for Federated Continual Learning}, in which a dual-encoder TransformerAE learns to map paired task weight vectors directly to merged-model weights---enabling zero-data, zero-gradient model composition. From 56 completed Stage-1 experiments across 39 curated loss functions and three data-overlap regimes, we establish several concrete findings:

\subsection{Confirmed Findings}

\textbf{FCL weight prediction is universally viable.} Every loss family tested achieves $\geq$75\% CNN ID accuracy after 5-epoch fine-tuning from predicted weights, a +66 percentage-point improvement over near-random initialization. This validates weight-space prediction as a practical FCL primitive.

\textbf{Loss choice is overlap-dependent.} No single loss dominates across all three overlap regimes. Log-norm regularization excels at low overlap (disjoint tasks); spectral FFT+MelSpec leads at partial overlap; optimal transport (LW\_Sinkhorn) wins at high overlap with the best overall accuracy (85.1\%) and lowest GW distance (13.4). Overlap-aware loss selection is therefore warranted in deployment.

\textbf{OT losses produce geometrically closer weights.} Sinkhorn-family losses reduce Gromov--Wasserstein distance to ground-truth weights by 47--79\% compared to MSE/MAE, indicating distributional shape preservation beyond pointwise fidelity. This geometric advantage translates to more robust performance at high overlap.

\textbf{Layerwise normalization is consistently beneficial.} All six LW variants match or exceed their global counterparts, with LWLN achieving the best average across completed runs (81.3\% vs.\ 80.6\% for MSE).

\subsection{Open Questions \textcolor{gray}{[to be addressed]}}

\begin{itemize}
    \item \textbf{Do topological losses outperform baselines?} DiffPers, RTD, and PersLandscape regularization are hypothesized to produce weight manifolds with topological structure more aligned with task structure. Results pending.
    \item \textbf{Does the winner profile hold at larger model sizes?} The multi-stage tournament (small $\to$ medium $\to$ large) will reveal whether overlap-dependent winners are robust or model-capacity artifacts.
    \item \textbf{Do HMMR-segmented training phases correlate with final accuracy?} Phase transition timing may predict convergence quality earlier than epoch 200.
\end{itemize}

\subsection{Limitations}

Our base CNN ($d=2{,}464$, MNIST) is intentionally small to enable large-scale loss comparison. Scaling to larger architectures requires resolving the quadratic complexity of Sinkhorn and the $O(n^3)$ persistence computation for DiffPers/RTD at larger $d$. The tournament protocol addresses this via progressive model-size expansion.

\subsection{Concluding Remarks}

This work makes the case that \emph{how} we measure error in weight space is not a secondary design choice---it determines the geometry of the learned prediction manifold, which in turn determines transfer efficiency across data distributions. The interplay between loss geometry, Gromov--Wasserstein structure, and downstream fine-tuning accuracy revealed here provides a foundation for principled loss design in federated and continual learning settings.

{
    \small
    \bibliographystyle{ieeenat_fullname}
    \bibliography{main}
}

% WARNING: do not forget to delete the supplementary pages from your submission 
% \input{sec/X_suppl}

\end{document}
